\subsubsection{Physically consistent scattering-kernel model}
\label{sec:beta-scattering-methods}

A single-band scattering fit cannot separate a scattering timescale from its
frequency scaling. If the broadening time is written phenomenologically as
\begin{equation}
\tau(\nu)=\tau_{1\,\mathrm{GHz}}
\left(\frac{\nu}{1\,\mathrm{GHz}}\right)^{-\alpha},
\end{equation}
then one observing band constrains only the value of $\tau$ across that band,
not the scaling index $\alpha$ independently. The CHIME/FRB--DSA-110
co-detections supply a $\sim$1\,GHz lever arm and therefore break this empirical
degeneracy. The physical interpretation of the scaling, however, requires one
more constraint: the pulse-broadening function (PBF) shape and $\alpha$ are not
independent parameters.

Both quantities follow from the same scattering physics. For an electron-density
fluctuation spectrum
\begin{equation}
P_n(q)\propto q^{-\beta},
\end{equation}
the inertial-range phase structure function scales as
\begin{equation}
D_\phi(\rho)\propto \rho^{\beta-2}.
\end{equation}
The transverse mutual coherence is then
$\Gamma(b)\simeq \exp[-D_\phi(b)/2]$; its Fourier transform gives the angular
brightness distribution of the scattered image, and the geometric delay
$t\propto\theta^2$ maps that angular distribution into the time-domain PBF. The
same $D_\phi$ also sets the diffractive scale and hence the frequency scaling of
the scattering time. For a thin screen in the inertial-range regime,
\begin{equation}
\alpha(\beta)=\frac{2\beta}{\beta-2},
\end{equation}
so Kolmogorov turbulence ($\beta=11/3$) gives
$\alpha=22/5\simeq4.4$, while the square-law limit ($\beta=4$) gives
$\alpha=4$ and a one-sided exponential PBF.

We therefore do not treat $\alpha$ as a free parameter inside a fixed PBF
family. A fixed exponential PBF already assumes square-law phase statistics and
therefore a narrow physical scaling near $\alpha=4$; allowing arbitrary
$\alpha$ values inside that kernel combines incompatible turbulence assumptions.
In particular, if the true PBF is heavier-tailed than an exponential, an
exponential fit can absorb the shape mismatch into a biased,
frequency-dependent $\widehat{\tau}$ and hence a biased apparent
$\widehat{\alpha}$ \citep{Cordes2025}.

Our primary joint model is therefore parameterized by the scattering strength and
the underlying scattering family, not by an independent PBF and independent
$\alpha$. In the simplest thin-screen implementation the shared sightline
parameters are
\begin{equation}
(\tau_{1\,\mathrm{GHz}},\beta),
\end{equation}
with the PBF shape and $\alpha(\beta)$ derived self-consistently at each
likelihood evaluation. The remaining pulse parameters are allowed to differ by
telescope: each band has its own arrival time $t_0$, intrinsic width $\zeta$,
residual dispersion offset $\delta\mathrm{DM}$, and component amplitudes. This
model is fit simultaneously to the CHIME/FRB and DSA-110 dynamic spectra.

This thin-screen mapping is the fiducial physical closure used here. Departures
from it---extended media, multiple screens, anisotropy, refraction, finite inner
scale, or a transition out of the inertial range---change the allowed relation
between PBF shape and frequency scaling. We treat such departures as model
limitations rather than allowing an unconstrained $\alpha$ to compensate for
them. Thus a formally shallow empirical scaling, especially one recovered under
a fixed exponential kernel, is not interpreted by itself as a turbulence spectral
index; it is evidence that the assumed scattering family or burst morphology is
incomplete.

\subsubsection{Gain-marginalized dynamic-spectrum likelihood}
\label{sec:beta-gainmarg}

The fit is performed on the two-dimensional dynamic spectra rather than on
band-averaged profiles. The model intensity in each frequency channel is a
scattered temporal profile multiplied by an unknown channel gain. This gain
absorbs the intrinsic burst spectrum, bandpass residuals, and diffractive
scintillation structure, none of which should be forced into the scattering
kernel.

We marginalize the per-channel gain analytically with a Gaussian prior. The
resulting per-channel evidence is a matched-filter likelihood for the temporal
shape after the unknown spectral amplitude has been integrated out. This reduces
the sampled dimensionality and prevents unresolved spectral structure from
appearing as correlated residuals in the time-domain scattering fit. We sample
the posterior with nested sampling as implemented in \texttt{dynesty}, and
quality-gate each fit using the reduced $\chi^2$ and lag-1 residual
autocorrelation; $R^2$ and residual normality are retained as diagnostic rather
than decisive gates.

\subsubsection{Sub-band validation of the frequency scaling}
\label{sec:beta-subband}

As an empirical cross-check, we also measure the scattering time independently in
frequency sub-bands. We split CHIME/FRB into four sub-bands and DSA-110 into
three sub-bands, collapse each sub-band to a one-dimensional profile, and fit a
scattered pulse with $\tau_i$ free in each sub-band. These sub-band fits do not
impose a common power law and do not share the full two-dimensional likelihood
used by the joint model.

The slope of the measured $\tau_i(\nu_i)$ points is therefore a diagnostic of
the observed frequency scaling. Agreement with the joint model supports the
recovered scaling; disagreement flags either profile evolution, insufficient
time resolution, low signal-to-noise, or a mismatch between the adopted
scattering family and the data. We use this sub-band reconstruction only as a
validation diagnostic, not as the primary estimator, because the CHIME tail can
be unresolved in faint or heavily scattered bursts and the DSA scattering time
can fall below the instrumental resolution at the top of the band.

\subsubsection{Multiple temporal components}
\label{sec:beta-multicomp}

Some FRBs contain multiple temporal components. If a component is omitted, the
fit can absorb it into the scattering tail, biasing both the recovered PBF shape
and the inferred frequency scaling. This is especially dangerous in a cross-band
fit: the apparent scaling can be driven by mismatched component structure
between CHIME/FRB and DSA-110 rather than by propagation physics.

We therefore extend the likelihood to allow $N$ temporal components per band.
Each component has its own arrival time and intrinsic width within that band,
while the scattering parameters describing the sightline are shared across all
components and both telescopes. Component counts are selected by Bayesian
evidence on a matched likelihood grid, so that one-component and multi-component
models have comparable evidence normalizations. Component arrival times within a
band are ordered and separated by a data-derived minimum spacing to keep the fit
identifiable.

When the evidence favors an additional component and the residuals whiten, we
interpret the component as burst morphology rather than scattering structure. If
the turbulence parameters rail or become multimodal in the one-component model
but unrail after the additional component is included, the single-component
result is rejected as a morphology-induced scattering bias.
